\documentclass{article}
\usepackage{amsmath}
\usepackage{amsthm}

\newtheorem{theorem}{Problem}

\begin{document}

\begin{theorem}
Compute $17^{-1}\pmod{83}$. Express your answer as a residue from $0$ to $82$, inclusive. (You may find it helpful to consider the fact that $17\cdot 5=85$.)
\end{theorem}

\begin{proof}
Step 0: Let's use the fact that $17\cdot 5=85$ to find the multiplicative inverse of 17 mod 83.
Step 1: The multiplicative inverse of 17 mod 83 is the number x such that $17x\equiv 1\pmod{83}$.
Step 2: And from the fact that $17\cdot 5=85$, we know that $17\cdot 5\equiv 2\pmod{83}$.
Step 3: Let's multiply both sides by $42$. We get $17\cdot 5\cdot42\equiv 2\cdot42\equiv1\pmod{83}$.
Step 4: So $x=5\cdot 42=210$ is a solution.
Step 5: But since we're looking for a number from 0 to 82, we need to find the remainder when 210 is divided by 83.
Step 6: That remainder is $210-83\cdot2=210-166=23$.
Step 7: So the multiplicative inverse of 17 mod 83 is 23.
Step 8: Right. And since $17\cdot 23=391$, we know that $17\cdot 23\equiv 391\equiv 1\pmod{83}$.
\end{proof}

\end{document}
